% Einheit 1 von 2 – Jahreszins, Monats- und Tageszins (bank/zinsrechnung/e1.jsonl, zusammenbau v0.1)
%% TODO Zweigzeile Teil 1: was hier gelernt wird (Satz in Schülersprache aus dem Einheitstitel) – Einheit 1
\einheitenkopf[e1]{Einheit 1 von 2 · Jahreszins, Monats- und Tageszins}
\zweigzeile{ab Kl. 8 (am Gymnasium ab Kl. 7) · P10}
\setcounter{aufgabe}{10}

%% TODO \verfahren{…}: Verfahrensüberschrift in Sachsprache fehlt in der Bank (2.3 b)
%% TODO Titel: Ich-kann-Satz fehlt in der Bank (Platzhalter: Kettenname) – Nr. 11
%% TODO Anweisung: ein Satz über der ersten Teilaufgabe fehlt in der Bank – Nr. 11
\begin{aufgabe}{Zinsen berechnen}
\begin{teile}
\teil Lena legt 850\,€ für ein Jahr auf ein Sparbuch, die Bank zahlt 2\,\% Zinsen. Wie viel Euro Zinsen bekommt sie? Was ist gesucht? Kreuze an, rechne nicht. \\ \kreuz{das Kapital (Grundwert)} \\ \kreuz{der Zinssatz (Prozentsatz)} \\ \kreuz{die Zinsen (Prozentwert)}
\end{teile}
%% TODO Darstellung für schwach fehlt in der Bank (zinsrechnung-e1-k1-s1-v1; Abschnitt „Für schwache Schüler“)
\swz{Sparbuch: 300\,€ zu 2\,\% – wie viel Euro Zinsen nach einem Jahr?}{}
%% TODO Darstellung für schwach fehlt in der Bank (zinsrechnung-e1-k1-s1-v2; Abschnitt „Für schwache Schüler“)
\swz{Sparbuch: 700\,€ zu 3\,\% – wie viel Euro Zinsen nach einem Jahr?}{}
%% TODO Darstellung für schwach fehlt in der Bank (zinsrechnung-e1-k1-s1-v3; Abschnitt „Für schwache Schüler“)
\swz{Sparbuch: 2\,000\,€ zu 4\,\% – wie viel Euro Zinsen nach einem Jahr?}{}
%% TODO Darstellung für schwach fehlt in der Bank (zinsrechnung-e1-k1-s1-v4; Abschnitt „Für schwache Schüler“)
\swz{Sparbuch: 900\,€ zu 2\,\% – wie viel Euro Zinsen nach einem Jahr?}{}
%% TODO Darstellung für schwach fehlt in der Bank (zinsrechnung-e1-k1-s1-v5; Abschnitt „Für schwache Schüler“)
\swz{Sparbuch: 5\,000\,€ zu 3\,\% – wie viel Euro Zinsen nach einem Jahr?}{}
\swfrage{Was bleibt gleich, was ändert sich?}
%% TODO Darstellung für schwach fehlt in der Bank (zinsrechnung-e1-k1-s2-v1; Abschnitt „Für schwache Schüler“)
\swz{Rechne über 1\,\%: 600\,€ zu 4\,\% – wie viel Euro Zinsen im Jahr?}{}
%% TODO Darstellung für schwach fehlt in der Bank (zinsrechnung-e1-k1-s3-v1; Abschnitt „Für schwache Schüler“)
\swz{Jana legt 700\,€ für ein Jahr zu 2\,\% an. Wie viel Euro hat Jana nach einem Jahr auf dem Konto?}{}
%% TODO Darstellung für schwach fehlt in der Bank (zinsrechnung-e1-k1-s4-v1; Abschnitt „Für schwache Schüler“)
\swz{900\,€ Kapital bringen in einem Jahr 27\,€ Zinsen. Wie hoch ist der Zinssatz?}{}
%% TODO Darstellung für schwach fehlt in der Bank (zinsrechnung-e1-k1-s5-v1; Abschnitt „Für schwache Schüler“)
\swz{Guthaben 4\,000\,€, Zinsen in einem Jahr 60\,€ – welcher Zinssatz?}{}
%% TODO Darstellung für schwach fehlt in der Bank (zinsrechnung-e1-k1-s6-v1; Abschnitt „Für schwache Schüler“)
\swz{21\,€ Zinsen im Jahr sind 3\,\% vom Kapital. Wie viel Euro Kapital?}{}
\swa{Weise mit der Tabelle nach, dass die Bank 3\,\% Zinsen im Jahr zahlt.}{\sachtabelle{lrrr}{Buchung & Einzahlung & Zinsen & Guthaben}{Jahresbeginn & 300,00\,€ & – & 300,00\,€\\ Jahresende & – & 9,00\,€ & 309,00\,€}}
%% TODO Darstellung für schwach fehlt in der Bank (zinsrechnung-e1-k1-s8-v1; Abschnitt „Für schwache Schüler“)
\swz{600\,€ zu 2\,\%, die Zinsen werden jedes Jahr ausgezahlt. Wie viel Euro Zinsen in 3 Jahren?}{}
%% TODO Darstellung für schwach fehlt in der Bank (zinsrechnung-e1-k1-s9-v1; Abschnitt „Für schwache Schüler“)
\swz{Für ein Motorroller leiht sich Familie Roth 1\,800\,€ für ein Jahr zu 6\,\%. Wie viel Euro muss sie nach einem Jahr zurückzahlen?}{}
%% TODO Darstellung für schwach fehlt in der Bank (zinsrechnung-e1-k1-s10-v1; Abschnitt „Für schwache Schüler“)
\swz{3\,000\,€ für ein Jahr: Bank A zahlt 2\,\%, Bank B 1,5\,\% und 20\,€ Startbonus. Welche Bank bringt mehr?}{}
\swz[3]{Ein Sparbetrag von 2\,600\,€ wird ein Jahr lang mit 2\,\% pro Jahr verzinst. Wie viel Euro Zinsen gibt es nach einem Jahr? \hfill (P10 2015 OS)}{}
\end{aufgabe}

%% TODO \verfahren{…}: Verfahrensüberschrift in Sachsprache fehlt in der Bank (2.3 b)
%% TODO Titel: Ich-kann-Satz fehlt in der Bank (Platzhalter: Kettenname) – Nr. 12
%% TODO Anweisung: ein Satz über der ersten Teilaufgabe fehlt in der Bank – Nr. 12
\begin{aufgabe}{Monats- und Tageszinsen}
\begin{teile}
\teil Geld liegt von Anfang März bis Ende August auf dem Konto. Welche Laufzeit? Kreuze an, rechne nicht. \\ \kreuz{ein Jahr – nicht teilen} \\ \kreuz{Monate – durch zwölf teilen} \\ \kreuz{Tage – durch dreihundertsechzig teilen}
\end{teile}
%% TODO Darstellung für schwach fehlt in der Bank (zinsrechnung-e1-k2-s1-v1; Abschnitt „Für schwache Schüler“)
\swz{2\,400\,€ zu 2\,\% für 5 Monate – wie viel Euro Zinsen?}{}
%% TODO Darstellung für schwach fehlt in der Bank (zinsrechnung-e1-k2-s1-v2; Abschnitt „Für schwache Schüler“)
\swz{1\,200\,€ zu 3\,\% für 4 Monate – wie viel Euro Zinsen?}{}
%% TODO Darstellung für schwach fehlt in der Bank (zinsrechnung-e1-k2-s1-v3; Abschnitt „Für schwache Schüler“)
\swz{3\,600\,€ zu 2\,\% für 7 Monate – wie viel Euro Zinsen?}{}
%% TODO Darstellung für schwach fehlt in der Bank (zinsrechnung-e1-k2-s1-v4; Abschnitt „Für schwache Schüler“)
\swz{6\,000\,€ zu 4\,\% für 3 Monate – wie viel Euro Zinsen?}{}
%% TODO Darstellung für schwach fehlt in der Bank (zinsrechnung-e1-k2-s1-v5; Abschnitt „Für schwache Schüler“)
\swz{4\,800\,€ zu 3\,\% für 9 Monate – wie viel Euro Zinsen?}{}
\swfrage{Was bleibt gleich, was ändert sich?}
%% TODO Darstellung für schwach fehlt in der Bank (zinsrechnung-e1-k2-s2-v1; Abschnitt „Für schwache Schüler“)
\swz{7\,200\,€ zu 2\,\% für 40 Tage, Bankjahr – wie viel Euro Zinsen?}{}
%% TODO Darstellung für schwach fehlt in der Bank (zinsrechnung-e1-k2-s3-v1; Abschnitt „Für schwache Schüler“)
\swz{5\,400\,€ zu 4\,\% vom 11. März bis zum 26. Mai, jeder Monat mit dreißig Tagen – wie viel Euro Zinsen?}{}
%% TODO Darstellung für schwach fehlt in der Bank (zinsrechnung-e1-k2-s4-v1; Abschnitt „Für schwache Schüler“)
\swz{3\,000\,€ bringen in 4 Monaten 20\,€ Zinsen. Welcher Zinssatz?}{}
\swz[3]{Nora hat 4\,200\,€ aus ihrem Nebenjob gespart. Sie legt das Geld von Anfang Mai bis Ende Oktober auf ein Tagesgeldkonto mit 2,5\,\% Zinsen im Jahr. Wie viel Euro Zinsen bekommt sie?}{}
\end{aufgabe}

%% TODO Titel: Ich-kann-Satz fehlt in der Bank (Platzhalter: Kettenname) – Nr. 13
%% TODO Anweisung: ein Satz über der ersten Teilaufgabe fehlt in der Bank – Nr. 13
\begin{aufgabe}{Überschlag: 1 \% von … , dann grob mal p}
%% TODO Darstellung für schwach fehlt in der Bank (zinsrechnung-e1-k3-s1-v1; Abschnitt „Für schwache Schüler“)
\swz{Überschlage: 3\,980\,€ zu 2\,\% – etwa wie viel Euro Zinsen im Jahr?}{}
\end{aufgabe}

%% TODO Titel: Ich-kann-Satz fehlt in der Bank (Platzhalter: Kettenname) – Nr. 14
%% TODO Anweisung: ein Satz über der ersten Teilaufgabe fehlt in der Bank – Nr. 14
\begin{aufgabe}{Zinsen berechnen – Fehler finden, Begründen, Anwendung}
\swz[3]{950\,€ zu 4\,\% – wie viel Euro Zinsen nach einem Jahr? Aylin rechnet: \rechnung{950 \cdot 0{,}04 &= 38 \\ 950 + 38 &= 988} Antwort: 988\,€. Finde den Fehler.}{}
\swfrage{Warum bringt 1\,\% von 20\,000\,€ mehr Zinsen als 4\,\% von 600\,€? Begründe ohne genau zu rechnen.}
\swz[3]{Frau Demir leiht sich 4\,000\,€ für ein Jahr zu 6\,\%. Sie kann jeden Monat 350\,€ zurücklegen. Reicht das Gesparte nach einem Jahr für die Rückzahlung?}{}
\end{aufgabe}

\uebersichtskasten{Zinsen: Zinsen sind der Prozentwert vom Kapital – das Kapital ist der Grundwert, der Zinssatz der Prozentsatz, gerechnet wird wie in der Prozentrechnung: 500 € zu 3 \% → Z = 500 · 0,03 = 15 € im Jahr. \\ Zinssatz: Zinsen geteilt durch Kapital, dann als Prozent schreiben: 45 € Zinsen von 1500 € → 45 : 1500 = 0,03 = 3 \%. \\ Guthaben nach einem Jahr: Kapital plus Zinsen – die Zinsen allein sind nicht das Guthaben: 500 € + 15 € = 515 €. \\ Kürzer als ein Jahr: Jahreszinsen durch 12 mal Monate, durch 360 mal Tage – die Bank rechnet das Jahr mit 360 Tagen und jeden Monat mit 30: 15 € Jahreszins → 4 Monate: 15 : 12 · 4 = 5 €.}
